% part: intuitionistic-logic
% chapter: propositions-as-types
% section: introduction

\documentclass[../../../include/open-logic-section]{subfiles}

\begin{document}

\olfileid{int}{pty}{int}

\olsection{تمهيد في مقابلة القضايا بالأنماط}

حساب لامبدا نموذج بسيط للحوسبة. وهو يعمل على حدود مبنية من المتغيرات. فإذا
كان~${\OLId{latin-looped}{N}}$ و${\OLId{latin-looped}{M}}$ حدين، كان~$(NM)$ حدًا كذلك (``تطبيق~${\OLId{latin-looped}{N}}$ على~${\OLId{latin-looped}{M}}$''). وإذا كان
${\OLId{latin-looped}{N}}$ حدًا، كان~$\lambd[{\OLId{latin-ordinary}{x}}][{\OLId{latin-looped}{N}}]$ حدًا. وإذا احتوى~${\OLId{latin-looped}{N}}$ المتغير~${\OLId{latin-ordinary}{x}}$، فإن وقائع
${\OLId{latin-ordinary}{x}}$ المقابلة في~$\lambd[{\OLId{latin-ordinary}{x}}][{\OLId{latin-looped}{N}}]$ تكون مقيّدة بالمعامل~$\lambd[{\OLId{latin-ordinary}{x}}]$، فلا تعود
حرة. ويقال إن حدًا من الصورة~$(\lambd[{\OLId{latin-ordinary}{x}}][{\OLId{latin-looped}{N}}]{\OLId{latin-looped}{M}})$ \emph{ينقبض بيتا} إلى
$\Subst{{\OLId{latin-looped}{N}}}{{\OLId{latin-looped}{M}}}{{\OLId{latin-ordinary}{x}}}$، أي إلى نتيجة استبدال~${\OLId{latin-looped}{M}}$ بكل وقوع حر لـ~${\OLId{latin-ordinary}{x}}$ في~${\OLId{latin-looped}{N}}$.
ويتحول~${\OLId{latin-looped}{N}}$ بتحويل بيتا إلى~${\OLId{latin-looped}{N}}'$ إذا نتج~${\OLId{latin-looped}{N}}'$ من~${\OLId{latin-looped}{N}}$ بانقباض بيتا لحد جزئي،
ونكتب~${\OLId{latin-looped}{N}}\redone {\OLId{latin-looped}{N}}'$. والحوسبة في حساب~$\lambd$ متتالية
${\OLId{latin-looped}{N}}\redone {\OLId{latin-looped}{N}}_{\OLMathNumeral{1}}\redone {\OLId{latin-looped}{N}}_{\OLMathNumeral{2}}\redone\dots$. فإذا انتهت المتتالية بحد~${\OLId{latin-looped}{N}}_{\OLId{latin-ordinary}{k}}$ لا
يمكن انقباض بيتا لأي حد جزئي منه، سُمّي \emph{سويًا} أو صورة سوية لـ~${\OLId{latin-looped}{N}}$.
وفكر في الحوسبة في حساب~$\lambd$ بوصفها متتاليات كهذه. تتوقف الحوسبة إذا
نتجت صورة سوية، وتكون تلك الصورة ناتج الحوسبة. ويتبين أن هذا النموذج البسيط
تام تورنغ: فكل دالة جزئية عودية يمكن حسابها بحد في حساب~$\lambd$.

اكتشف هاسكل كاري ووليام هوارد، كل على حدة، تشابهًا وثيقًا بين الاستنباط
الطبيعي وحساب~$\lambd$. فالحد~$\lambd[{\OLId{latin-ordinary}{x}}][{\OLId{latin-looped}{N}}]$ دالة بحسب الحدس: المتغير~${\OLId{latin-ordinary}{x}}$
هو المعطى، و${\OLId{latin-looped}{N}}$ يبين كيف نحسب القيمة انطلاقًا منه. ويكون
$(\lambd[{\OLId{latin-ordinary}{x}}][{\OLId{latin-looped}{N}}]{\OLId{latin-looped}{M}})$ تطبيق الدالة~$\lambd[{\OLId{latin-ordinary}{x}}][{\OLId{latin-looped}{N}}]$ على المعطى~${\OLId{latin-looped}{M}}$، ويبين
انقباض بيتا كيفية تقييمه: نستبدل~${\OLId{latin-looped}{M}}$ بـ~${\OLId{latin-ordinary}{x}}$ في~${\OLId{latin-looped}{N}}$، ثم نواصل تقييم الحد
الناتج. فكر الآن في هذه الدوال بوصفها ذات مجال ومدى ثابتين. فإذا كان مجال
$\lambd[{\OLId{latin-ordinary}{x}}][{\OLId{latin-looped}{N}}]$ هو~${\OLId{latin-looped}{A}}$ ومداه~${\OLId{latin-looped}{B}}$، وجب أن نعد~${\OLId{latin-ordinary}{x}}$ متغيرًا لا يأخذ إلا
معطيات من~${\OLId{latin-looped}{A}}$، وأن نعد~${\OLId{latin-looped}{N}}$ منتجًا قيمًا من~${\OLId{latin-looped}{B}}$. ولذلك لا يكون
$(\lambd[{\OLId{latin-ordinary}{x}}][{\OLId{latin-looped}{N}}]{\OLId{latin-looped}{M}})$ ذا معنى إلا إذا كان~$\lambd[{\OLId{latin-ordinary}{x}}][{\OLId{latin-looped}{N}}]$ دالة~${\OLId{latin-looped}{A}}\to {\OLId{latin-looped}{B}}$ وكان
${\OLId{latin-looped}{M}}\in {\OLId{latin-looped}{A}}$. وإذا أضفنا هذه المعلومات إلى حساب~$\lambd$ حصلنا على حساب
$\lambd$ \emph{المنمّط}. وليست كل عبارة~$(\lambd[{\OLId{latin-ordinary}{x}}][{\OLId{latin-looped}{N}}]{\OLId{latin-looped}{M}})$ مشروعة في
الحساب المنمّط، بل تقتصر المشروعة على العبارات التي تتلاءم فيها ``أنماط''
$\lambd[{\OLId{latin-ordinary}{x}}][{\OLId{latin-looped}{N}}]$ و${\OLId{latin-looped}{M}}$؛ أي يكون نمط~$\lambd[{\OLId{latin-ordinary}{x}}][{\OLId{latin-looped}{N}}]$ هو~${\OLId{latin-looped}{A}}\to {\OLId{latin-looped}{B}}$ ونمط~${\OLId{latin-looped}{M}}$
هو~${\OLId{latin-looped}{A}}$. ولا يكون نمط~$\lambd[{\OLId{latin-ordinary}{x}}][{\OLId{latin-looped}{N}}]$ هو~${\OLId{latin-looped}{A}}\to {\OLId{latin-looped}{B}}$ إلا إذا كان نمط~${\OLId{latin-ordinary}{x}}$ هو
${\OLId{latin-looped}{A}}$ ونمط~${\OLId{latin-looped}{N}}$ هو~${\OLId{latin-looped}{B}}$. وبوجه عام ليست كل عبارة~$({\OLId{latin-looped}{M}}_{\OLMathNumeral{1}}{\OLId{latin-looped}{M}}_{\OLMathNumeral{2}})$ مشروعة، بل
تقتصر الحدود المشروعة منها على ما كان فيه نمط~${\OLId{latin-looped}{M}}_{\OLMathNumeral{1}}$ هو~${\OLId{latin-looped}{A}}\to {\OLId{latin-looped}{B}}$ ونمط
${\OLId{latin-looped}{M}}_{\OLMathNumeral{2}}$ هو~${\OLId{latin-looped}{A}}$؛ وعندئذ يكون نمط~$({\OLId{latin-looped}{M}}_{\OLMathNumeral{1}}{\OLId{latin-looped}{M}}_{\OLMathNumeral{2}})$ هو~${\OLId{latin-looped}{B}}$.

تقابل قواعد التنميط هذه بوضوح~!!{derivation}s في الاستنباط الطبيعي، حيث
تمثّل الأنماط~!!{formula}s، وتقابل~!!{derivation}s نفسها حدود~$\lambd$.
فمثلًا، يقابل~!!a{derivation} لـ~$!A\lif !B$ حدًا
$\lambd[\typeof{{\OLId{latin-ordinary}{x}}}{!A}][{\OLId{latin-looped}{N}}]$ نمطه الدالي~$!A\lif !B$. وتقابل قواعد
الاستدلال في الاستنباط الطبيعي قواعد التنميط في حساب لامبدا المنمّط، مثلًا:
\begin{prooftree}
  \AxiomC{$\Discharge{!A}{{\OLId{latin-ordinary}{x}}}$}
  \DeduceC{$!B$}
  \DischargeRule{\Intro\lif}{x}
  \UnaryInfC{$!A \lif !B$}
  \DisplayProof\bottomAlignProof
  \quad\text{تقابل}\quad
  \Axiom${\OLId{latin-ordinary}{x}}: !A \fCenter {\OLId{latin-looped}{N}}: !B$
  \UnaryInf$ \fCenter \lambd[\typeof{{\OLId{latin-ordinary}{x}}}{!A}][{\OLId{latin-looped}{N}}] : !A \lif !B$
\end{prooftree}
وتعني القاعدة على اليمين أنه إذا كان نمط~${\OLId{latin-ordinary}{x}}$ هو~$!A$ ونمط~${\OLId{latin-looped}{N}}$ هو~$!B$،
كان نمط~$\lambd[\typeof{{\OLId{latin-ordinary}{x}}}{!A}][{\OLId{latin-looped}{N}}]$ هو~$!A\lif !B$.

تقابل قاعدة~$\Elim\lif$ قاعدة التنميط لحدود التطبيق، أي:
\[
  \AxiomC{$!A \lif !B$}
  \AxiomC{$!A$}
  \RightLabel{\Elim\lif}
  \BinaryInfC{$!B$}
  \DisplayProof
  \quad\text{تقابل}\quad
  \Axiom$\fCenter {\OLId{latin-looped}{M}}_{\OLMathNumeral{1}}: !A \lif !B$
  \Axiom$\fCenter {\OLId{latin-looped}{M}}_{\OLMathNumeral{2}}: !A$
  \BinaryInf$ \fCenter ({\OLId{latin-looped}{M}}_{\OLMathNumeral{1}}{\OLId{latin-looped}{M}}_{\OLMathNumeral{2}}) : !B$
  \DisplayProof
\]
إذا تبعت قاعدةُ~$\Elim\lif{}$ قاعدةَ~$\Intro\lif{}$ فورًا، أمكن تبسيط
!!{derivation}:
\begin{gather*}
  \AxiomC{$\Discharge{!A}{{\OLId{latin-ordinary}{x}}}$}
  \DeduceC{$!B$}
  \DischargeRule{\Intro\lif}{{\OLId{latin-ordinary}{x}}}
  \UnaryInfC{$!A \lif !B$}
  \AxiomC{}
  \DeduceC{$!A$}
  \RightLabel{\Elim\lif}
  \BinaryInfC{$!B$}
  \DisplayProof
  \quad\redone\quad
  \AxiomC{}
  \DeduceC{$!A$}
  \DeduceC{$!B$}
  \DisplayProof
\end{gather*}
وهذا يقابل انقباض بيتا لحدود لامبدا:
\[
(\lambd[\typeof{{\OLId{latin-ordinary}{x}}}{!A}][{\OLId{latin-looped}{N}}]{\OLId{latin-looped}{M}}) \quad\redone\quad
\Subst{{\OLId{latin-looped}{N}}}{{\OLId{latin-looped}{M}}}{{\OLId{latin-ordinary}{x}}}.
\]

وتوجد مقابلات مماثلة بين قواعد~$\land$ والأنماط ``الضربية''، وبين قواعد
$\lor$ والأنماط ``الجمعية''.

تسمى هذه المقابلة بين حدود حساب لامبدا البسيط المنمّط و!!{derivation}s
الاستنباط الطبيعي مقابلة ``كاري--هوارد'' أو ``البراهين بوصفها برامج'' أو
``القضايا بوصفها أنماطًا''. وإلى جانب مقابلة~!!{formula}s (القضايا) للأنماط
والبراهين للحدود، يمكن تلخيص المقابلات هكذا:
\begin{center}
  \begin{tabular}{c c}
    المنطق & البرنامج \\
    \hline
    القضية/!!{formula} & النمط \\
    البرهان & الحد \\
    الفرض & المتغير \\
    الفرض المفرّغ & المتغير المقيّد \\
    الفرض غير المفرّغ & المتغير الحر \\
    $!A \lif !B$ & النمط الدالي \\
    $!A \land !B$ & النمط الضربي \\
    $!A \lor !B$ & النمط الجمعي \\
    $\lfalse$ & النمط الخالي \\
    \hline
  \end{tabular}
\end{center}

مقابلة كاري--هوارد إحدى ركائز مساعدات البرهان الآلية وفاحصات أنماط البرامج،
لأن فحص برهان يشهد لقضية، كما فعلنا أعلاه، لا يعدو فحص ما إذا كان برنامج
(حد) يحمل النمط المعلن.
\end{document}
